\documentclass{situationist-preprint}

\title{Natural Beauty and Scientific Cognitivism}
\msid{AES-11-a}
\msfamily{AES-11}
\msdiscipline{Aesthetics}
\mstype{Principal manuscript}
\mskeywords{natural beauty; environmental aesthetics; scientific cognitivism; representation; categories of art; qualified observers}

\begin{document}
\maketitle

\begin{abstract}
Is nature beautiful in the way art is? We argue that it is not, and that it differs in two respects. First, natural beauty is not representational: natural objects represent nothing, because representation requires a maker whose intentions, realized in the object, give it content, and natural objects have no maker. Resemblances in nature, such as a cloud that looks like a face, are occasions for seeing-in but not representations. Second, the appreciation of natural beauty requires scientific rather than art-historical knowledge. On the response-dependent account of beauty, an object is beautiful if it is disposed to please qualified observers who perceive it in its correct category. For artworks, correct categories are fixed by facts of art history; for natural objects, they are fixed by the facts of natural history that the sciences describe. A qualified observer of nature therefore needs scientific knowledge, as a qualified observer of art needs art-historical knowledge. We defend this scientific cognitivism against the arousal model, against the view that nature has no correct categories, and against formalism.
\end{abstract}

\section{The problem}

Kant held that natural beauty has a priority over the beauty of art: a person who takes an immediate interest in the beauty of nature shows a disposition favorable to moral feeling, which interest in art does not show \citep[§42]{Kant1790}. Hegel reversed the ranking, holding that the beauty of art, as born of the spirit, stands higher than the beauty of nature (Hegel, \emph{Lectures on Aesthetics}, Introduction), and analytic aesthetics, for most of the twentieth century, followed Hegel in practice by attending almost exclusively to art. \citet{Hepburn1966} called attention to this neglect and argued that the aesthetic appreciation of nature has features the art-centered model misses: nature is unframed, the appreciator is inside it rather than in front of it, and its objects are not designed to be looked at.

The question we address is whether natural beauty is beauty of the same kind as the beauty of art, appreciated in the same way. Two answers have shaped the subsequent debate. \citet{Carlson1979,Carlson2000} defended the \emph{natural environmental model}, often called scientific cognitivism: as the appreciation of art requires knowledge of art history, so the appropriate appreciation of nature requires knowledge of natural history, as provided by the natural sciences. Non-cognitivists have rejected this requirement. \citet{NCarroll1993} argued that we may be appropriately moved by nature, for example by the grandeur of a waterfall, through an emotional arousal that requires only common-sense knowledge; \citet{Brady2003} defended a role for imagination and perceptual engagement free of scientific concepts; and \citet{Budd2002} argued that the appreciation of nature is free in a way the appreciation of art is not, because natural objects have no correct category of the kind that art-historical facts fix for works. Others have connected scientific cognitivism to the thesis that unspoiled nature is always beautiful \citep{Parsons2002} and to the requirement that nature be appreciated on its own terms \citep{Saito1998}.

We defend scientific cognitivism, and we derive it from the account of beauty in AES-02 together with the role of categories in aesthetic perception. We add an independent claim that the debate has largely taken for granted and that settles the other half of our question: unlike much art, nature is not representational.

\section{Framework}

\begin{definition}[Representation]\label{d:rep}
An object \emph{represents} $p$ if it has $p$ as content in virtue of a maker's successful intention that it have that content, realized in the object's features under conventions or capacities of recognition available to its intended audience.
\end{definition}

Definition~\ref{d:rep} applies to representation in art the account of meaning defended in AES-09, on which intentions fix content when they succeed. It requires a maker; it does not require that the content be conventional, since depictive content is grounded in our capacities of recognition rather than in conventions alone.

\begin{definition}[Seeing-in]\label{d:seein}
\emph{Seeing-in} is the experience of seeing something in a marked surface or other perceptible array, as a face is seen in a cloud or a figure in a stain on a wall. It does not require that the array represent what is seen in it.
\end{definition}

The distinction is Wollheim's: seeing-in is a perceptual capacity prior to representation, and representation occurs when a standard of correctness, set by the intentions of a maker, determines what is to be seen in the array \citep{Wollheim1980}.

\begin{principle}[Response-dependence]\label{p:rd}
An object is beautiful if and only if it is disposed to give aesthetic pleasure to qualified observers who perceive it, in suitable conditions, in its correct category (AES-02).
\end{principle}

\begin{principle}[Categories]\label{p:cat}
The aesthetic properties an object has depend on the category in which it is correctly perceived: the same perceptible features may be standard in one category and contra-standard in another, and their aesthetic character differs accordingly \citep{Walton1970}.
\end{principle}

\begin{principle}[Correctness of categories]\label{p:fix}
The correct category of an object is the category it belongs to in virtue of the facts that explain its having the features it has. For artworks these are facts about its making and its place in an artistic tradition; for natural objects, facts about the natural processes that formed it.
\end{principle}

Principle~\ref{p:cat} is Walton's thesis that a work seen as a Cubist painting and the same work seen in a category of its own would have different aesthetic properties, and that the category in which the work was made, and in which its maker intended it to be perceived, is the correct one. Principle~\ref{p:fix} generalizes the ground of correctness. What makes a category correct for a work, in Walton's account, is a set of historical facts: the maker's intentions, the categories recognized in the society in which it was produced, and the features it has that are standard in that category. The underlying idea is that the correct category is the one that reflects why the object has the features it does. For natural objects, the facts that explain why they have their features are facts of geology, biology and physics.

\section{The argument}

\begin{lemma}\label{l:norep}
Natural objects are not representational.
\end{lemma}

\begin{proof}
By Definition~\ref{d:rep}, an object represents something only in virtue of a maker's successful intention realized in its features. A natural object, such as a mountain, a cloud or a shell, is not the product of any agent's intentional act, and so has no maker in the relevant sense. It therefore has no representational content. Where we see a face in a cloud or a crouching lion in a rock formation, what occurs is seeing-in (Definition~\ref{d:seein}), which requires only a resemblance that triggers our recognitional capacities; there is no standard of correctness that makes the face what is to be seen. Another observer who sees a ship in the same cloud is not mistaken.
\end{proof}

Lemma~\ref{l:norep} marks a first difference between natural and artistic beauty. Much of the beauty of art is the beauty of a representation: of how a subject is presented, of the relation between what is shown and how it is shown, and of the achievement of showing it. A landscape painting is beautiful partly in virtue of how it represents the landscape. None of this is available in the landscape itself. Not all art is representational, so Lemma~\ref{l:norep} does not by itself separate nature from all of art; it separates natural beauty from the representational beauty that a large part of art possesses and that no natural object can have.

\begin{lemma}\label{l:cat}
The correct categories of natural objects are fixed by natural history, as described by the natural sciences.
\end{lemma}

\begin{proof}
By Principle~\ref{p:fix}, an object's correct category is the one it belongs to in virtue of the facts that explain its features. For a natural object, these are the processes that formed it: a whale is shaped as it is because it is an air-breathing mammal adapted to the sea, a glacial valley because of glacial erosion, a coral reef because of the growth of colonial organisms. The categories that sort objects by these processes are those of the natural sciences: species, landforms, ecosystems and their kinds. Art-historical categories, such as genres and styles, sort objects by their place in a tradition of making, and no natural object has such a place.
\end{proof}

\begin{theorem}\label{t:main}
Nature is not beautiful in the way art is. Natural beauty is not representational, and the appreciation of natural beauty requires scientific rather than art-historical knowledge.
\end{theorem}

\begin{proof}
The first clause is Lemma~\ref{l:norep}. For the second, by Principle~\ref{p:rd} whether a natural object is beautiful depends on the responses of qualified observers who perceive it in its correct category. By Principle~\ref{p:cat}, what the observer perceives, and so what she responds to, depends on the category in which she perceives it, and an observer who perceives the object in an incorrect category responds to aesthetic properties it does not have. By Lemma~\ref{l:cat}, the correct categories of natural objects are scientific categories. An observer can perceive an object in a category only if she grasps the category and can recognize the object as belonging to it. The qualified observer of nature therefore needs knowledge of the relevant scientific categories, and art-historical knowledge, which supplies categories natural objects do not belong to, does not meet this need.
\end{proof}

Carlson's own example illustrates the second clause. A whale perceived as a fish appears sluggish, bulky and ungainly, since its features are contra-standard for fish; perceived as a mammal, the same features are standard and the animal appears graceful and powerful \citep{Carlson2000}. On the account of AES-02, the second perception is the one whose responses determine whether the whale is beautiful, because it is the perception of a qualified observer in the correct category. The knowledge required is not encyclopedic. It is the knowledge needed to perceive the object as the kind of thing it is, and its demands scale with the object: little is needed to see that a tree is a tree, more to see a landform as the product of glaciation.

\section{Objections and replies}

\begin{objection}[Arousal]\label{o:arousal}
We can be moved by the grandeur of a waterfall without knowing anything about the hydrology of waterfalls, and our being moved is appropriate. Scientific knowledge is not required for appropriate appreciation \citep{NCarroll1993}.
\end{objection}

\begin{reply}
The responses Carroll describes are appropriate, and we do not deny them. A person with no knowledge of hydrology can be moved by a waterfall's force and scale, and those properties may be perceptible in the category \enquote{waterfall}, which common sense supplies. The claim of Theorem~\ref{t:main} concerns the appreciation of natural beauty in the sense of AES-02, which is fixed by qualified observers. Common-sense categories are often correct as far as they go, and the knowledge they embody is rudimentary natural history. Where they diverge from scientific categories, as in the case of the whale, the common-sense response is mistaken, and it is mistaken for the reason scientific cognitivism gives.
\end{reply}

\begin{objection}[Freedom]\label{o:freedom}
Natural objects have no correct category, because no one made them to be perceived in one. The appreciator of nature is free to perceive a natural object under any category she likes \citep{Budd2002}.
\end{objection}

\begin{reply}
The objection identifies the ground of correctness with a maker's intention. Principle~\ref{p:fix} identifies it with the facts that explain an object's features, of which a maker's intention is one case. The objection also has a consequence that its proponents do not accept: if there were no correct category, there would be no fact about whether the whale is graceful or ungainly, and the judgment that it is graceful would be as correct as its denial. Our practice does not treat the difference between the informed and the uninformed observer of nature in this way. The freedom that nature allows is real but lies elsewhere: in the choice of what to attend to and from where, since nature is unframed \citep{Hepburn1966}.
\end{reply}

\begin{objection}[Formalism]\label{o:formal}
Many natural beauties, such as the color of a sunset or the symmetry of a snowflake, are formal and can be appreciated without any knowledge of the object's kind. The Kantian tradition of free beauty describes exactly this appreciation.
\end{objection}

\begin{reply}
Some aesthetic properties of natural objects may be category-independent, and to that extent the appreciation of them requires no knowledge. Kant's account of free beauty is examined in the companion manuscript AES-11-b. But the formalist cannot generalize the point: the gracefulness of the whale, the delicacy of an alpine plant in its habitat, and the drama of a geological formation depend on the category in which they are perceived. The formal properties that are category-independent are a small part of natural beauty, and a natural aesthetic confined to them would be impoverished in the way an aesthetic of painting confined to color harmony would be.
\end{reply}

\begin{objection}[Other knowledge]\label{o:folk}
Indigenous and local traditions often have their own categories for natural objects, in narrative and practical terms. Scientific cognitivism privileges one tradition over others \citep{Saito1998}.
\end{objection}

\begin{reply}
Theorem~\ref{t:main} requires knowledge of the categories that explain the features of natural objects. Traditional natural knowledge often supplies such categories, sometimes in great detail, and where it does it counts as natural history in the sense required. What the theorem excludes is knowledge of the wrong kind: categories that sort natural objects by the stories told about them or by their resemblance to artworks. These may enrich the experience of a place, but they are not the categories in which its objects are correctly perceived.
\end{reply}

\section{Conclusion}

Natural beauty differs from the beauty of art in two respects. It is not representational, because natural objects have no maker whose successful intentions could give them content, and resemblances in nature are occasions for seeing-in, not representations. And its appreciation requires scientific rather than art-historical knowledge, because beauty is fixed by qualified observers perceiving objects in their correct categories, and the correct categories of natural objects are fixed by the processes that formed them. Nature is beautiful, and it is beautiful on its own terms.

\begin{thebibliography}{99}\small
\bibitem[Brady(2003)]{Brady2003} Brady, E. (2003). \emph{Aesthetics of the Natural Environment}. Edinburgh University Press.
\bibitem[Budd(2002)]{Budd2002} Budd, M. (2002). \emph{The Aesthetic Appreciation of Nature: Essays on the Aesthetics of Nature}. Clarendon Press.
\bibitem[Carlson(1979)]{Carlson1979} Carlson, A. (1979). Appreciation and the natural environment. \emph{The Journal of Aesthetics and Art Criticism}, 37(3), 267--275.
\bibitem[Carlson(2000)]{Carlson2000} Carlson, A. (2000). \emph{Aesthetics and the Environment: The Appreciation of Nature, Art and Architecture}. Routledge.
\bibitem[Carroll(1993)]{NCarroll1993} Carroll, N. (1993). On being moved by nature: Between religion and natural history. In S. Kemal and I. Gaskell (Eds.), \emph{Landscape, Natural Beauty and the Arts}. Cambridge University Press.
\bibitem[Hepburn(1966)]{Hepburn1966} Hepburn, R.~W. (1966). Contemporary aesthetics and the neglect of natural beauty. In B. Williams and A. Montefiore (Eds.), \emph{British Analytical Philosophy}. Routledge \& Kegan Paul.
\bibitem[Kant(1790)]{Kant1790} Kant, I. (1790). \emph{Critique of the Power of Judgment}. Trans. P. Guyer and E. Matthews. Cambridge University Press, 2000.
\bibitem[Parsons(2002)]{Parsons2002} Parsons, G. (2002). Nature appreciation, science, and positive aesthetics. \emph{The British Journal of Aesthetics}, 42(3), 279--295.
\bibitem[Saito(1998)]{Saito1998} Saito, Y. (1998). Appreciating nature on its own terms. \emph{Environmental Ethics}, 20(2), 135--149.
\bibitem[Walton(1970)]{Walton1970} Walton, K.~L. (1970). Categories of art. \emph{The Philosophical Review}, 79(3), 334--367.
\bibitem[Wollheim(1980)]{Wollheim1980} Wollheim, R. (1980). Seeing-as, seeing-in, and pictorial representation. In \emph{Art and Its Objects}, 2nd ed. Cambridge University Press.
\end{thebibliography}

\end{document}
